\SourceLeaf{375}%
\[
\begin{gathered}
\xi=\frac{rx}{aa},\qquad\eta=\frac{ry}{bb},\qquad\zeta=\frac{rz}{cc},\\[4pt]
rr=aa\xi\xi+bb\eta\eta+cc\zeta\zeta,
\end{gathered}
\]
also
\[
ds=\frac{aabbcc\,d\sigma}{{(aa\xi\xi+bb\eta\eta+cc\zeta\zeta}^{2}}.
\]
\needspace{14\baselineskip}

Soll die ganze Oberfläche des Sphäroids bestimmt werden, so ist
\[
d\sigma=\cos\varphi\,.\,d\varphi\,.\,d\lambda
\]
zu setzen und die doppelte Integration von $\lambda=0$ bis $\lambda=360^0$ und von $\varphi=-90^0$ bis $\varphi=+90^0$ auszuführen. Die erstere ist leicht und man hat dann, $\sin\varphi=x$ gesetzt, die ganze Oberfläche
\[
=\int\frac{\pi aabbcc\,.(aa+bb-(aa+bb-2cc)xx)\,dx}{(aa-(aa-cc)xx)^{\frac32}\,.(bb-(bb-cc)xx)^{\frac32}}
\]
von $x=-1$ bis $x=+1$.

Dieses Integral verwandelt sich leicht in folgendes (allgemein)
\[
\begin{aligned}
&\pi aabbcc\,\frac{\left(\dfrac1{aa}+\dfrac1{bb}-\dfrac1{cc}\right)x+\left(\dfrac{(aa-cc)(bb-cc)}{aabbcc}\right)x^3}{\BareRoot{\{(aa-(aa-cc)xx)(bb-(bb-cc)xx)\}}}\\[7pt]
&\quad+\pi aabbcc\int\frac{\dfrac1{cc}-\dfrac{(aa-cc)(bb-cc)}{aabbcc}xx}{\BareRoot{\{(aa-(aa-cc)xx)(bb-(bb-cc)xx)\}}}\,dx,
\end{aligned}
\]
also von $x=-1$ bis $x=+1$ in folgenden Werth
\[
2\pi cc+2\pi\int\frac{aabb-(aa-cc)(bb-cc)xx}{\BareRoot{\{(aa-(aa-cc)xx)(bb-(bb-cc)xx)\}}}\,dx,
\]
das Integral von $x=0$ bis $x=1$ genommen.
